\documentclass[a4paper,12pt]{article}
% add all packages
\usepackage[a4paper,bottom=35mm]{geometry}
\usepackage{amsmath}
\usepackage{amsthm}
\usepackage{mdframed}
\usepackage{textcomp} 
\usepackage{siunitx}          		
\usepackage[parfill]{parskip}    		
\usepackage{graphicx}
\usepackage{xcolor}
\usepackage{caption}
\usepackage{bbm}	
\usepackage{amssymb}
\usepackage{paralist}
\usepackage{color}
\usepackage[utf8]{inputenc}
\usepackage[english,german]{babel}
\usepackage{mathtools}
\usepackage{nicefrac}
\usepackage{wrapfig}
\usepackage[mathscr]{euscript}
\usepackage{ytableau}
\usepackage{mathtools}
\usepackage{fancyhdr}
\usepackage{url}
\usepackage{hyperref}

% definiert die Sprache
\selectlanguage{German}
\setlocalecaption{german}{figure}{Figur}
\setlocalecaption{german}{table}{Tabelle}

% add command for different dates in Swiss fashion
\newcommand{\monthwordlong}[1]{\ifcase#1 \or January\or February\or March\or April\or May\or June\or July\or August\or September\or October\or November\or December\fi}
\newcommand{\monthwordshort}[1]{\ifcase#1\or Jan\or Feb\or Mar\or Apr\or May\or Jun\or Jul\or Aug\or Sept\or Okt\or Nov\or Dez\fi} 
\newcommand{\datelong}{\number\day. \monthwordlong{\the\month} \number\year}
\newcommand{\dateshort}{\number\day. \monthwordshort{\the\month} \number\year}

% capital roman numbers
\newcommand{\uproman}[1]{\uppercase\expandafter{\romannumeral#1}}

% define color of links inside a text
\definecolor{myblue}{RGB}{0, 0, 80}
\definecolor{plotblue}{RGB}{100, 100, 255}
\hypersetup{
	colorlinks=true, 
	linkcolor=myblue,
	filecolor=magenta,      
	urlcolor=plotblue,
	citecolor=myblue,
	pdfpagemode=FullScreen,
}

% placing a figure at the top of clear page
\makeatletter
\setlength{\@fptop}{0pt}
\makeatother

% add single blankpage
\def\blankpage{\clearpage\thispagestyle{empty}\addtocounter{page}{-1}\null\clearpage}

% Lückentext erstellen
\newcommand{\lineortext}[1]{\hspace{0,1cm}\underline{\hspace{0,25cm}\hphantom{#1}\hspace{0,25cm}}\hspace{0,1cm}}
% \renewcommand{\lineortext}[1]{\hspace{0,1cm}\underline{\hspace{0,25cm}#1\hspace{0,25cm}}\hspace{0,1cm}}

% math commands
\theoremstyle{definition} % don't write everything in italic
\newtheorem{myfact}{Fact}
\newtheorem{mydef}{Definition}[section]
\newenvironment{fdef}{\begin{mdframed}\begin{mydef}}{\end{mydef}\end{mdframed}}
\newtheorem{mytheo}{Theorem}
\newenvironment{ftheo}{\begin{mdframed}\begin{mytheo}}{\end{mytheo}\end{mdframed}}
\newtheorem{myprop}{Proposition}
\newtheorem{mylemma}{Lemma}[section]
\newtheorem{myaux}{Auxiliary Lemma}
\newtheorem{mycor}{Corollary}
\newtheorem{myre}{Remark}[section]
\newtheorem{myex}{Example}
\newtheorem{mynot}{Notation}

\newtheorem{satz}{Satz}[section]
\newenvironment{fsatz}{\begin{mdframed}\begin{satz}}{\end{satz}\end{mdframed}}
\newtheorem{aufg}{Aufgabe}
\newtheorem{bsp}{Beispiel}

\newcommand{\R}{\mathbb{R}}   
\newcommand{\N}{\mathbb{N}}
\newcommand{\Z}{\mathbb{Z}}
\newcommand{\Q}{\mathbb{Q}}
\newcommand{\C}{\mathbb{C}}
\newcommand{\Lim}[1]{\raisebox{0.5ex}{\scalebox{0.8}{$\displaystyle \lim_{#1}\;$}}}
\newcommand{\vvec}[2]{\left(\begin{array}{c} #1 \\ #2 \end{array}\right)}               % two dimensional vector
\newcommand{\vvector}[3]{\left(\begin{array}{c} #1 \\ #2 \\ #3 \end{array}\right)}      % three dimensional vector
\newcommand{\irow}[1]{\begin{smallmatrix}(#1)\end{smallmatrix}}                         % inline row vector
\newcommand{\icol}[1]{\left(\begin{smallmatrix}#1\end{smallmatrix}\right)}              % inline column vector
\DeclarePairedDelimiter\ceil{\lceil}{\rceil}                                            % ceil- function
\DeclarePairedDelimiter\floor{\lfloor}{\rfloor}                                         % floor - function


% header and footers
\pagestyle{fancy}
\fancyhf{}
\fancyhead[L]{}
% \fancyhead[C]{header}
\fancyhead[R]{}
\fancyfoot[R]{\thepage}
% check for more information: https://de.overleaf.com/learn/latex/Headers_and_footers


% AdminCommands
\newcommand{\name}{Jonas Guler}
\newcommand{\examdate}{13. Sept 2022}
\newcommand{\topic}{Vektorgeometrie Probeprüfung}
\newcommand{\subject}{Mathematik}
\newcommand{\class}{KSH - 3GaLaWa}
\newcommand{\school}{KSH}

\begin{document}
\thispagestyle{fancy}
\lhead{\sf \large \topic \\ \small \name, \examdate}
\rhead{\sf \class \\  Name: \hspace*{4cm}}
\begin{small}
	\begin{center}
		\boxed{\textbf{Wichtige Informationen:}}  \\
		\vspace{2pt}
		Der Rechenweg muss vollständig und nachvollziehbar sein.\\
		Die Schlussresultate sind auf zwei Nachkommastellen zu runden oder vollständig gekürzt anzugeben.\\
		Lösen Sie die Prüfung mit einen schwarzen oder blauen Kugelschreiber. Resultate in Bleistift, roter oder grüner Farbe werden nicht berücksichtigt.\\
		Schreiben Sie auf \textbf{jedes} Blatt (inklusive Aufgabenblatt) Ihren Namen.\\
		Es sind keine Hilfsmittel erlaubt.\\
		
		! VIEL ERFOLG !
	\end{center}
\end{small}

\Aufgabe{1 (2+4)}
\begin{itemize}
	\item[a)] Betrachte die folgende Vektorgleichung.
	\[ \frac{1}{2}(\vec{a}+\vec{c}) - \frac{3}{4}(3\vec{b}-\vec{a}) = \frac{1}{4}(4\vec{b}+\vec{c}) + \vec{a} \]
	Drücke den Vektor $\vec{a}$ in Abhängigkeit von $\vec{b}$ und $\vec{c}$ aus.
	
	\item[b)] Beweise vektoriell, dass ein Viereck, in dem sich die Diagonalen halbieren, ein Parallelogramm ist.
\end{itemize}

\newpage
\pagestyle{fancy}
\lhead{\sf \large \topic}
\rhead{\sf \class}


\Aufgabe{2 (2)}
\begin{itemize}
	\item[a)] Gegeben sind die Vektoren $\vec{a} = \vvector{0}{-2}{8}, \vec{b} = \vvector{-1}{3}{2}$ und $\vec{c} = \vvector{3}{-5}{-3}$.\\ Berechne den Vektor $\vec{v} = \vec{a}-2\vec{b}+3\vec{c}$.
	
%	\item[b)] Weiter sei der Vektor $\vec{w} = \vvector{x}{y}{\nicefrac{5}{2}}$ gegeben. Berechne $x$ und $y$, sodass $\vec{v}$ und $\vec{w}$ kollinear sind.
\end{itemize}

\newpage
\Aufgabe{3 (6)}
\begin{itemize}
	\item[a)] Löse folgendes Gleichungssystem nach $x$, $y$ und $z$ auf.
	\begin{eqnarray*}
        2x-3y & = & -6\\
        1x+5y+5z & = & 2\\
        -3x+7y-4z & = & 5 
    \end{eqnarray*}
	%\[ \vec{a} = \vvector{2}{1}{-3}, \vec{b} = \vvector{-3}{5}{7}, \vec{c} = \vvector{0}{5}{-4}, \vec{d} = \vvector{-6}{2}{5} \]
	
	%\item[b)] Gegeben sind die Punkte $A = (4,4,3)$ und $B = (2,0,-1)$. In welchen Punkten durchstösst die x-Achse die Kugel mit Durchmesser $AB = |\vec{AB}|$.
\end{itemize}

\end{document}
